\section{正交与正交补：把空间分成互不干扰的方向}
\label{sec:03-Linear-Algebra/03-Orthogonality/01}

前面全部讨论都围绕一件事：哪些向量可以通过线性组合产生哪些向量。向量空间告诉你什么能加、什么不能加，基告诉你最少需要几个方向来描述整个空间。但有一个重要的问题一直没问：两个方向有多接近？一个向量有多长？近似的质量用一个什么数字来衡量？

没有长度和角度的线性代数，好比只知道路名和路与路之间如何连接，却不知道每条路到底有多长。你从 A 到 B 可以走很多条路径，但不知道哪条最短。本单元为向量空间加上几何尺度。所有新概念——正交、投影、最小二乘——都从点积这一个运算生长出来。

\subsection{点积：用一个数字衡量两个方向的亲密程度}

在 \(\R^n\) 中，对
\[
x=(x_1,\ldots,x_n)^\trans,\qquad
y=(y_1,\ldots,y_n)^\trans,
\]
定义
\[
\langle x,y\rangle=x^\trans y
=x_1y_1+\cdots+x_ny_n.
\]
这个简单的运算同时给出长度和角度。取 \(x=y\)：
\[
\|x\|_2^2=x^\trans x,\qquad
\|x\|_2=\sqrt{x^\trans x}.
\]
这就是欧氏长度。两点之间的距离随之确定为 \(\|x-y\|_2\)。

方向之间的关系交给余弦。当 \(x,y\) 都非零时，定义
\[
\cos\theta=\frac{x^\trans y}{\|x\|_2\|y\|_2}.
\]
分母把长度的影响去掉，剩下的值只反映方向。为什么这个分式一定落在 \([-1,1]\) 之间？因为 Cauchy--Schwarz 不等式
\[
|x^\trans y|\le\|x\|_2\|y\|_2
\]
保证了分子不会超过分母。

这个不等式不需要另外承认。对任意实数 \(t\)，向量 \(x-ty\) 的长度平方天然非负：
\[
0\le\|x-ty\|_2^2
=\|x\|_2^2-2t\,x^\trans y+t^2\|y\|_2^2.
\]
右端是关于 \(t\) 的二次函数。它恒非负意味着图像永远不会跑到横轴下方，因此判别式不能是正的：
\[
(2x^\trans y)^2-4\|y\|_2^2\|x\|_2^2\le0,
\]
整理就是 \((x^\trans y)^2\le\|x\|_2^2\|y\|_2^2\)。两边开方即得。若 \(y=0\)，两端都是零，不等式显然成立。

Cauchy--Schwarz 的推导过程自身就解释了它的含义：如果有一个向量 \(y\)，沿着 \(y\) 的方向去近似 \(x\)，残差的平方不可能小于零——这个残差的下界恰好由 \(x^\trans y\) 和 \(\|y\|_2\) 决定。Cauchy--Schwarz 等于在说：投影之前先检查一下你最多能解释多少。

本节专注欧氏点积。更一般的内积可以替换长度和垂直的含义（例如函数空间里用积分作内积），但只要满足线性、对称和正定性三条，后面的所有推理结构保持不变。

\subsection{垂直等同于点积为零}

我们从初中就知道直角三角形满足 \(a^2+b^2=c^2\)。拿到向量语境：如果 \(x\) 和 \(y\) 构成直角，那么把两者首尾相接得到的 \(x+y\)，其长度的平方应当等于两部分长度平方之和：
\[
\|x+y\|_2^2=\|x\|_2^2+\|y\|_2^2.
\]
展开左边：
\[
\begin{aligned}
\|x+y\|_2^2
&=(x+y)^\trans(x+y)\\
&=x^\trans x+2x^\trans y+y^\trans y.
\end{aligned}
\]
和勾股关系相比，中间多了一个 \(2x^\trans y\)。所以勾股关系等价于中间项为零：
\[
x^\trans y=0.
\]
正交的代数判据就是这么来的。它不是凭空定义的符号约定，而是``直角三角形的勾股关系对向量加法成立''的代数翻译。记作 \(x\perp y\)。

例如
\[
x=\begin{pmatrix}1\\2\\2\end{pmatrix},
\qquad
y=\begin{pmatrix}2\\-1\\0\end{pmatrix}
\]
满足 \(x^\trans y=2-2+0=0\)。验算长度：
\[
\|x+y\|_2^2=\|(3,1,2)^\trans\|_2^2=9+1+4=14,
\]
\[
\|x\|_2^2+\|y\|_2^2=9+5=14.
\]
勾股关系无一例外地成立。

一个细节不要忽略：零向量和任何向量都正交（点积都是零），但零向量没有方向，所以不能往夹角公式里塞。这个例外在后面讨论正交基的时候还会回来——正交基里的每个向量都必须是货真价实的非零方向。

\subsection{正交方向为什么比普通方向好用得多}

取非零向量 \(q_1,\ldots,q_k\)，两两正交。设
\[
c_1q_1+\cdots+c_kq_k=0.
\]
两边与 \(q_j\) 做点积。除了 \(c_j\|q_j\|_2^2\) 这一项，其他项全部因为正交而消失：
\[
c_j\|q_j\|_2^2=0.
\]
\(q_j\neq0\) 所以 \(c_j=0\)，对每个 \(j\) 都如此。结论：非零的两两正交向量自动线性无关。正交比线性无关更强——后者只说不能互相表示，前者还说它们在任何线性组合中的系数可以直接通过点积分离开来。

如果再让每个向量的长度都是 \(1\)——这就是正交归一向量。此时若
\[
v=c_1q_1+\cdots+c_kq_k,
\]
直接与 \(q_j\) 做点积：
\[
c_j=q_j^\trans v.
\]
在普通基下取坐标需要解一个线性方程组；在正交归一基下，每个坐标只需一次点积。计算代价从 \(O(n^3)\) 直接塌缩到 \(O(n^2)\)。后面几节所有构造正交基的努力，都源于这一个乘法复杂度的落差。

这个性质在数据中有一层更微妙的对应。数据表的两列可以看作 \(n\) 维向量。如果两列都已减去各自的样本均值，\(x^\trans y/n\) 恰好是样本协方差——衡量两列数据共同涨落程度的一个数字。于是 \(x^\trans y=0\) 表示这批样本中的线性共同变化为零。在几何上这意味着两个数据向量垂直；在统计上对应不相关。但需要清醒的是：这个正交关系依赖中心化方式和所选内积，它不能单独推出两个变量独立——除非再有联合 Gaussian 这样额外的结构绑定。几何的垂直落到数据里对应的是线性不相关，不是毫无关系。

\subsection{从正交向量到正交子空间}

单向量之间的垂直很容易检查。两个子空间之间的垂直要更谨慎。

设 \(S,T\subseteq\R^n\) 是两个子空间。只有对任意 \(s\in S\) 和任意 \(t\in T\) 都成立 \(s^\trans t=0\)，才称 \(S\) 与 \(T\) 正交。``任意''不能省。两个平面相交成一条直线，并不意味着两个平面正交；只要能在两个平面中各自找到一个向量、二者的点积非零，定义就被破坏。

最常见的错误是：看到两个子空间的基向量彼此正交，就宣布两个子空间正交。基向量正交并不等于子空间中所有向量两两正交——还需要确认基向量的任意线性组合之间是否仍然正交。只有当一个子空间中任一向量与另一子空间中任一向量都垂直，才有资格说两个子空间正交。

\subsection{正交补：补全缺少的所有垂直方向}

给定子空间 \(S\)，收集所有与 \(S\) 中所有向量都正交的向量：
\[
S^\perp=\{x\in\R^n:x^\trans s=0\text{ 对所有 }s\in S\}.
\]
\(S^\perp\) 自身也是一个子空间——零向量在内、对标量乘法和加法都封闭。它叫作 \(S\) 的正交补。

名字中的``补''不仅表示垂直，更重要的意思是：它补足了 \(S\) 缺少的、构成整个 \(\R^n\) 所需的全部方向。在有限维欧氏空间中：
\[
\R^n=S\oplus S^\perp.
\]
这条结论不需要额外假设，可以从已有工具直接推出。首先，如果 \(v\) 同时属于 \(S\) 和 \(S^\perp\)，那么 \(v^\trans v=0\)，即 \(v=0\)——两者只交于零向量。接着，取 \(S\) 的一组基作为行构成矩阵，这个矩阵的零空间恰好是 \(S^\perp\)。由秩--零度定理，\(\dim S+\dim S^\perp=n\)。只交于零向量、维数之和又恰好等于 \(n\)——两个子空间共同把 \(\R^n\) 完整地剖成不相重叠的两半。

于是每个 \(b\in\R^n\) 都能唯一地写成
\[
b=p+e,\qquad p\in S,\; e\in S^\perp.
\]
\(p\) 是 \(b\) 在 \(S\) 中的分量，\(e\) 是 \(b\) 沿 \(S\) 的垂直方向上的分量。刚才这一整段推导只说清楚了一种分解存在且唯一，还没有给出怎么算。下一节投影会把这个分解变成可执行的算法，并进一步告诉你 \(p\) 恰好是 \(S\) 中离 \(b\) 欧氏距离最近的点。

\subsection{矩阵的四个基本子空间如何配对}

现在可以把正交补和上一节的四个基本子空间对接起来。设 \(A\) 是 \(m\times n\) 实矩阵。

若 \(x\in\mathcal N(A)\)，则 \(Ax=0\)。把 \(A\) 的第 \(i\) 行记作 \(r_i^\trans\)，方程 \(Ax=0\) 的第 \(i\) 个分量就是
\[
r_i^\trans x=0.
\]
\(x\) 与 \(A\) 的每一行都正交，自然也正交于这些行的任意线性组合——后者恰好是整个行空间。因此
\[
\mathcal N(A)=\operatorname{Row}(A)^\perp.
\]

把同样的逻辑用到 \(A^\trans\) 上：
\[
\mathcal N(A^\trans)=\operatorname{Col}(A)^\perp.
\]

这四条正交关系把四个基本子空间锁成两对互为正交补的搭档：
\[
\R^n=\operatorname{Row}(A)\oplus\mathcal N(A),
\]
\[
\R^m=\operatorname{Col}(A)\oplus\mathcal N(A^\trans).
\]
第一组在输入空间 \(\R^n\)，第二组在输出空间 \(\R^m\)。初学者最容易犯的错误就是把两对混在一起：以为零空间的正交补是列空间，或者把 \(\mathcal N(A)\) 和 \(\mathcal N(A^\trans)\) 当成一对。它们分处两个环境空间——向量长度不同的空间不能直接比正交关系。

两对正交补揭露了 \(A\) 作为线性变换的完整收支表：

\begin{itemize}
\tightlist
\item
  输入空间 \(\R^n\)：行空间承载能传到输出端的有效信号，零空间承载被 \(A\) 完全抹除的信息。
\item
  输出空间 \(\R^m\)：列空间是 \(A\) 实际能到达的靶面，左零空间是无法被 \(A\) 命中的方向。
\end{itemize}

行空间和零空间互为正交补，解释了为什么秩--零度定理的维数能这样干净地闭合——两个子空间不仅在维数上互补，
在几何上也被强制分到了不共享任何非零向量的两个半空间里。

\subsection{一段完整分解：拆一个向量到一维子空间和它的正交补}

设 \(S=\operatorname{span}\{(1,1,0)^\trans\}\subseteq\R^3\)。要判定正交补：\(x=(x_1,x_2,x_3)^\trans\in S^\perp\) 当且仅当与生成向量的点积为零：
\[
x_1+x_2=0,
\]
所以 \(S^\perp\) 是一个二维平面，由条件 \(x_2=-x_1\) 刻画的全部向量构成。

现在把 \(b=(2,0,1)^\trans\) 拆成两个分量。设 \(p=t(1,1,0)^\trans\in S\)，残差
\[
e=b-p=(2-t,-t,1)^\trans
\]
必须落在 \(S^\perp\)。带入条件：
\[
(1,1,0)\,e=(2-t)+(-t)+0=2-2t=0,
\]
解出 \(t=1\)。因此
\[
b=\underbrace{(1,1,0)^\trans}_{p\in S}
+\underbrace{(1,-1,1)^\trans}_{e\in S^\perp}.
\]
两个分量分属正交的两个子空间，点积：\(1\cdot1+1\cdot(-1)+0\cdot1=0\)，验证通过。

这个例子只用了一条直线，靠解一个方程得到分解。下一节的投影会针对任意维数的子空间给出一个不需要凑答的通用方法，并证明这条唯一分解中，\(p\) 还是 \(S\) 中与 \(b\) 欧氏距离最近的向量。

\subsection{正交性在不同内积下有不同面貌}

本节从头到尾都在用 \(x^\trans y\) 作为内积，得到的是标准欧氏空间中的长度、角度和垂直。但如果问题环境变了——比如对不同的方向应该给予不同的重视程度，或者函数空间中的内积是积分——正交的定义必须跟着变。只要是合法的内积（双线性、对称、正定），正交补的存在和直和分解照样成立。但``哪个方向是垂直的''会跟着内积的权重一起调整。

这个区分在下一节的加权最小二乘和后续的数值分析中都会反复出现：投影永远是对某个内积的最近点。选对了内积，模型才可能接近真实；选错了内积，数学推导全部正确，结论却可能毫无用处。

延伸阅读：MIT 18.06 Lecture 14：Orthogonal Vectors and Subspaces。
